\pdfoutput=1
\documentclass[11pt]{article}
\usepackage[margin=1in]{geometry}
\usepackage{amsmath,amssymb,amsthm,mathtools}
\usepackage[hidelinks]{hyperref}
\hypersetup{pdftitle={A Sourced BRST Cut Charge and a Finite Graviton Transition on Compact Milne}}
\newtheorem{theorem}{Theorem}
\newtheorem{lemma}[theorem]{Lemma}
\newtheorem{proposition}[theorem]{Proposition}
\theoremstyle{remark}
\newtheorem{remark}[theorem]{Remark}
\newcommand{\dd}{\mathrm d}
\newcommand{\supp}{\operatorname{supp}}
\newcommand{\ad}{\operatorname{ad}}
\newcommand{\tr}{\operatorname{tr}}
\newcommand{\Ric}{\operatorname{Ric}}
\newcommand{\TT}{\mathrm{TT}}
\newcommand{\WF}{\operatorname{WF}}
\newcommand{\Cq}{Q_{\mathrm c}}
\newcommand{\Db}{D_{\mathrm{bulk}}}
\newcommand{\one}{\mathbf 1}
\newcommand{\Lie}{\mathcal L}
\title{A Sourced BRST Cut Charge\\and a Finite Graviton Transition on Compact Milne}
\author{}
\date{}
\begin{document}
\maketitle

\begin{abstract}
A source-dependent master identity, a conserved interior charge, and an observed transition are different objects in perturbative gravity. We compare them in a single causally renormalized pure Einstein theory. With the ordinary local constant-coupling quantum master normalization as input, a coherent coupling profile admits a local jet descent of its master defect. The past part of the actual defect corrects the free BRST charge. On the dynamical quotient the result equals the independently varied, smoothly smeared Noether charge and is nilpotent at the retained formal orders. On compact hyperbolic Milne we construct a compatible full graded Hadamard realization and three smooth, geometrically marked tensor packets for which finite unswitched Einstein evolution has a nonzero connected one-to-two coefficient. Smooth exterior samplers select a positive two-graviton contribution of order \(\hbar\kappa^2\). It uses one state and retains the cross interference when a bookkeeping cut is moved. The result is a finite perturbative test, not an asymptotic scattering rate, a separately quantized boundary theory, or a full elastic one-loop amplitude.
\end{abstract}

\section{Setup and claims}\label{sec:setup}

We use perturbative algebraic quantum field theory for four-dimensional pure Einstein gravity. The reference spacetime \(M\) is a regular globally hyperbolic vacuum slab with compact oriented Cauchy surface. The explicit experiment uses a subslab of compact hyperbolic Milne
\[
       g_0=-\dd t^2+t^2\gamma,\qquad
       \operatorname{sec}(\gamma)=-1,\qquad t>0 .
\]
The interaction is expanded in a canonically normalized metric perturbation with parameter \(\kappa\); the first Einstein interaction is cubic in the tensor sector. We use the gauge-fixed nonminimal Batalin--Vilkovisky (BV) splitting, including the free field BRST term, on a single Epstein--Glaser time-ordering prescription \(T\). Ghost, antighost, multiplier and antifield insertions are kept until the graded dynamical quotient. In particular \(Q_0=\delta_0+\gamma_0\), including the Koszul--Tate term \(\delta_0\). We keep the classical Einstein BV action fixed and put the loop normalizations in \(T\). We assume the usual causal, local covariant, real and field-equation normalizations, the free Ward compatibility of the covariances, and the ordinary local pure-gravity quantum master equation (QME) normalization on constant-coupling germs, in the sense of \cite[Sec.~3.4]{BFR}, in this same split. These are finite-order hypotheses on \(T\) and its local marked family, not claims to remove an anomaly for arbitrary matter content or globally on arbitrary off-shell sectors. All statements below are coefficientwise through the orders used: the once-interacting observable and its first nonzero selected norm-square, and the local Ward identities through the first loop and second coupling order. The strict square refers to its leading coefficient; no norm-convergent series is asserted.

A smooth compact temporal function \(f\) is \(1\) on an open plateau containing all preparations, samplers and bookkeeping cuts. It changes only in two disjoint exterior collars. Use the localized interaction
\[
           V_f=\sum_{n\geq1}\kappa^n V_n[f^n],
           \qquad R_f=R_{V_f},
\]
with the free quadratic operator unchanged. The locality tests and profiles needed in the marked Ward family use the same \(T\). No physical source or intervention is applied at an interior cut. A superscript \(-\) on an exterior insertion denotes only its past collar part, not a negative-frequency projection. We use left graded commutators \([A,B]_{\star}\) and \(D_q=(i/\hbar)[q,\cdot]_{\star}\), choosing the free charge \(q_0\) to implement \(\gamma_0\) after imposing the free equations. All on-shell statements below take this quotient, including removal of external antifields as dynamical symbols.

\begin{theorem}[Same-theory sourced charge]\label{thm:charge}
The full local master defect of \(V_f\) admits a local representative supported in the two switching collars. Write \(J_f\) for the negative of that defect, in the common causal prescription. On interior on-shell dynamical insertions the physical bulk differential is implemented by
\begin{equation}\label{eq:keycharge}
       \Cq=q_0-R_f(J_f^-).
\end{equation}
In this equation \(\Cq\) is also the retarded image of the independently varied, renormalized interior Einstein Noether charge, smeared by a compact unit temporal one-form in the plateau. More generally one may use closed compact one-forms of the same cut cohomology class there. It is the same on-shell element for all such smearings; \(\Cq\star\Cq=0\) at the retained orders. Equality is not claimed between off-shell Noether densities or between separately renormalized theories on two sides.
\end{theorem}

\begin{theorem}[Finite Milne experiment]\label{thm:experiment}
On compact hyperbolic Milne one can choose a compatible graded free Hadamard realization with positive physical cohomology, one associated formal physical preparation, and three smooth orthogonal tensor packet labels \(a,b,c\). There are disjoint, positive-width, smoothly sampled exterior times, both in the plateau, for which the connected, earlier-readout-subtracted \(a\)-to-\((b,c)\) coefficient of the unswitched Einstein dynamics has the form
\begin{equation}\label{eq:cubicform}
   -i\kappa\sqrt{\hbar}\,\mathfrak a,\qquad \mathfrak a\ne0 .
\end{equation}
The resulting selected two-graviton contribution at order \(\hbar\kappa^2\) is \(\hbar\kappa^2|\mathfrak a|^2\). All intermediate smooth artificial cut presentations give this same coefficient with the same physical pairing. Splitting the interval before taking the norm-square retains the interference between its two pieces.
\end{theorem}

\paragraph{Antecedents and distinction.}
Rejzner identifies the on-shell interacting BV differential as a commutator with the interacting BRST charge, including general relativity \cite{Rejzner,BFR}. Hollands and Taslimi Tehrani develop local Ward, renormalized current and switched-charge arguments in Yang--Mills theory \cite{Hollands,TT}; Rejzner--Schiavina give the smooth-boundary modified quantum master equation and the retarded past-source term \cite[Thm.~4.16, Lemma~5.1 and Sec.~5.4]{RS}. We use these identities rather than asserting the BRST commutator or the modified master equation as new. The additions here are their comparison in the same coherent Einstein prescription with explicit source support, the compatible graded Milne benchmark, and a nonzero measured Einstein tree with its physical interference in the common cut representation. A separate local Einstein--scalar benchmark treated free gauge transport and a classical matter subcocycle \cite{P19}; a free shear example compared positive covariance and coherent sewing \cite{P53}. Neither supplies the finite interacting Einstein process considered here.

\section{The sourced master identity and the interior charge}\label{sec:ward}

\subsection{A local current with the actual support}

In the sign convention of \cite[Lemma~4.8, Thm.~4.16]{RS}, the anomalous Master Ward identity for the one \(T\) gives
\begin{align}
 Q_0 S_T(V_f)
  &=\frac{i}{\hbar}\,T\bigl(e^{iV_f/\hbar},M_f\bigr),\nonumber\\
 M_f&=Q_0 V_f+\tfrac12\{V_f,V_f\}_{\rm std}
                  -i\hbar A_T(V_f).
 \label{eq:AWI}
\end{align}
Here braces denote the BV antibracket and \(A_T\) the local anomaly of this same extension of the time-ordered products. We use the coherent powers \(f^n\) because a constant-source germ \(f=a\) then equals the constant-coupling germ with coupling \(a\kappa\) at every homogeneous degree. This includes brackets and all local descendants with several vertices. Equal powers \(f\) on all distinct couplings would not have this property. Pure Einstein below means the full gauge-fixed BV Lagrangian, not only its graviton self-interaction; in particular the source jets of the antifield and ghost insertions are kept.

The classical part of the current is the negative of the first two terms of \(M_f\). This identifies its variational convention, not just its sign. In the generalized representative of \cite[Props.~4.3, 4.9]{RS}, the standard current differs from twice the source current by the contraction of the variational boundary potential (\(J_{\rm std}=2J_{\rm src}-\iota_Q\theta\), with that reference's test smearing). The nonminimal relative comparison also contains
\[
     B_0(V_f)=Q_0V_f-\{L_0,V_f\}_{\rm std}.
\]
Indeed its classically adjusted current
\(J_{\rm std}(L_0+V_f)-J_{\rm std}(L_0)-B_0(V_f)\)
is \(-Q_0V_f-\{V_f,V_f\}_{\rm std}/2\) by expansion of the bracket. We use this comparison classically; we do not apply its bare formula to quantum contact terms. The BV split, Euler contact terms and this boundary correction may not be dropped just because the profile is constant in the interior.

Here is an explicit descent of the full residual as a local functional of the switching profile. At a retained homogeneous coefficient of ghost number one, let \(m_f\) be a finite-jet local density for \(M_f\), including the coincident-vertex terms of \(A_T\). On the constant-jet subspace, \(m_f=f^n m_n\) at coupling degree \(n\). For any constant profile, the normalized unmarked QME holds as an integrated identity under compact off-shell field variations in an interior chart. Thus all Euler derivatives of \(m_n\) in the fluctuating BV fields are zero. Its field-independent part is zero by ghost number, with no external ghost constants. If \(\Theta_{m_n}\) is the boundary term in the first-variation formula for the density, scale the fluctuating fields \(\Phi\) by \(r\) and integrate:
\[
        m_n=\dd b_n,\qquad
        b_n(\Phi)=\int_0^1\Theta_{m_n}(r\Phi;\Phi)\,\dd r .
\]
A common covariant integration-by-parts convention uses the smooth background connection on the jet bundles. This does not assume that the density itself vanishes for constant \(f\).

Subtract \(\dd(f^n b_n)\) from \(m_f\). The difference vanishes when all positive-order source jets vanish. Write it as
\[
 P(f,\xi)-P(f,0)
     =\sum_{|I|\geq1}\xi_I\int_0^1
                      \partial_{\xi_I} P(f,r\xi)\,\dd r ,
       \qquad \xi_I=\nabla_I f,
\]
at the retained local finite-jet order. Integrate by parts on the chosen differentiated factor until one derivative of \(f\) remains; the other source jets in the coefficient stay in place. Thus in our convention
\begin{equation}
      J_f=-M_f=\int_M \dd f\wedge j[f] .
      \label{eq:descent}
\end{equation}
This is an off-shell identity of integrated local functionals; \(j[f]\) is a profile-dependent finite-jet three-form, not a universal profile-independent Noether current. No division by \(f\) or \(\dd f\) is involved. A marked derivative of \eqref{eq:descent} is taken in the same joint causal extension (the same \(T\) and Wick convention), not with each resulting graph separately finite-part renormalized. This is important for field-equation/propagator contacts. In particular \eqref{eq:descent} permits the genuine support split \(J_f=J_f^-+J_f^+\) of \cite[Remark~4.6, Lemma~5.1]{RS}.

For a buffered local insertion \(F\) supported inside the plateau, \cite[Eq.~(37)]{RS} then gives, before the dynamical quotient,
\begin{equation}
 Q_0 R_f(F)
   =R_f(\Db F)+\frac{i}{\hbar}[R_f(J_f^-),R_f(F)]_{\star},
 \qquad
 \Db F=Q_0F+\{V_f,F\}-i\hbar \Delta_{V_f}^{\rm ren}F .
 \label{eq:RS}
\end{equation}
The future term cancels by causal factorization; the past term reverses its order relative to time ordering, producing the commutator. Local coincident terms are retained in \(\Delta_{V_f}^{\rm ren}\). Equation~\eqref{eq:RS} for multilocal generators means the appropriate marked descendants of the same relative \(S\)-matrix, or the derivation on the algebra generated by the local on-shell insertions after the proof below. It is not an assumption that the ordinary pointwise product of interacting multilocal functionals obeys the same formula term by term.

\subsection{Noether insertion normalized inside the plateau}

Independently start with the gauge-fixed Einstein BRST Noether three-form \(j_{\rm N}\) from the variational boundary potential and divergence choice of \cite[App.~B]{BFR}. We use a compact bump one-form \(\eta_{\rm c}\) in a buffered collar with unit oriented integral along the time coordinate; compact spatial sections make this a compact current smear. Use a second, independent test \(\beta\) to mark the divergence within the plateau, not to turn off its interaction. The renormalized field-equation insertion and marked Ward relations in \cite[Eqs.~(40), (43), (67), (70)]{Rejzner} give the on-shell identity
\begin{equation}
    R_f\bigl((\dd j_{\rm N})(\beta)\bigr)
       = i\hbar R_f\bigl(A(V;\beta)\bigr),\qquad
    A(V;\beta)=\sum_{n}\Delta^{n}(V^{\otimes n};V(\beta)).
    \label{eq:markdiv}
\end{equation}
The notation in the marked slot means the infinitesimal local Lagrangian probe around the constant plateau; it is normalized there and is not a replacement of the coherent physical cutoff \(f\). All equalities can be understood on compact mark supports and their causal neighborhoods. We are \emph{not} imposing \(\Delta(V_f)=0\) on the switched presentation. The characteristic support of local anomalies means the coefficients in \eqref{eq:markdiv} depend only on the constant-coupling germ where the test or its derivatives are evaluated. Contributions outside the plateau belong to \eqref{eq:RS} instead.

To normalize the current insertion inside the plateau, first move derivatives off \(\beta\) in its local anomaly density. The anomaly is understood in the \emph{same nonminimal split} as \eqref{eq:AWI}, using the free-Ward-compatible time ordering (the tilded anomaly in \cite[Sec.~3.4, Eq.~(65)]{Rejzner}). A marked profile equal to one near the compactly supported variations of all BV fields reads the unmarked constant-coupling anomaly in that same family. The classical master equation holds separately for the fixed classical action, as an integrated germ identity. The constant-coupling QME normalized by \(T\) therefore implies that this anomaly density has vanishing Euler derivative \emph{off shell}. Source-gradient terms vanish only in this chosen interior germ. The field-scaling formula used above now gives a local three-form \(k_V\) with \(A(V;\beta)=(\dd k_V)(\beta)\) as integrated local functionals before the equation quotient. This is the jet version of eliminating the zeroth marked integral of a total-diagonal distribution; compare \cite[Sec.~4.6, Lemma~9]{Hollands}. The separate field-equation contacts in the actual current Ward identity are not discarded by this off-shell choice. Define \(j_{\rm Q}=j_{\rm N}-i\hbar k_V\). It is conserved in the common interacting on-shell algebra on the plateau, with the same prescription at every cut. It does not cancel the actual \(J_f\) in a region with profile gradients.

For completeness the charge implementation uses the \emph{differentiated} marked identity, not the false replacement \(\Delta_{V_f}^{\rm ren}F=0\). In the variation of \eqref{eq:markdiv} with a local test insertion \(F\) hold the primitive \(k_V\) at the physical \(V\) fixed in the independently specified current slot. The responses differentiating the retarded normalization of the undifferentiated \(A(V;\beta)\) and its subtraction \((\dd k_V)(\beta)\) cancel. The remaining local anomaly descendants are both
\begin{equation}
  \sum_n n\,\Delta^n(V^{\otimes(n-1)},F;V(\beta)),
   \qquad
  \sum_n\Delta^n(V^{\otimes n};F(\beta)),
  \label{eq:descendants}
\end{equation}
in addition to classical bracket and Euler terms. These are respectively the variation in an interaction slot and in the probe slot of the local anomaly. The retarded-minus-advanced (GLZ) identity takes two compact past/future mark profiles with difference \(1\) on a neighborhood of \(\supp F\) and derivative zero there. On the coincident distributions in \eqref{eq:descendants}, it replaces the marked terms precisely by the unmarked linearization \(\Delta_{V_f}^{\rm ren}F\), as in \cite[Eqs.~(43), (70)]{Rejzner}. All Euler terms are removed using the renormalized field-equation insertion identity after differentiation. This argument is local on the plateau and retains the two descendants, which is why a global zero-anomaly hypothesis on the switched \(V_f\) is unnecessary. The classical bracket contribution is fixed by the variational Noether identity; the resulting charge satisfies
\begin{equation}
 \Cq^{\rm N}=R_f\left(\int\eta_{\rm c}\wedge j_{\rm Q}\right),
 \qquad
 [R_f(\Db F)]_{\rm os}
  =\frac{i}{\hbar}[\Cq^{\rm N},R_f(F)]_{\star,{\rm os}} .
 \label{eq:inner}
\end{equation}
The conversion of the printed right bracket includes its degree: for an odd charge, \([Q,F]_\star=-(-1)^{|F|}[F,Q]_\star\). All operations here use the left convention fixed for \(Q_0\), not a degree-independent sign replacement. The commutator comparison with a switched defect parallels \cite[Thms.~15--17, Eq.~(102)]{TT}; we have specified its locality, variational reference and switched source here in a pure Einstein chart.

\subsection{On-shell equality, nilpotence, and smooth recutting}

On antifield-independent dynamical representatives the free equations remove \(\delta_0\) on the left of \eqref{eq:RS}; the remaining \(\gamma_0\) is the free charge commutator. Comparing \eqref{eq:RS} with \eqref{eq:inner} therefore gives
\begin{equation}\label{eq:center}
 [\Cq^{\rm N}-q_0+R_f(J_f^-),R_f(F)]_{\star,{\rm os}}=0 .
\end{equation}
Use the \emph{full gauge-fixed graded dynamical algebra} at this step, not only its gauge-invariant subalgebra. Algebraic multipliers are eliminated by their equations in the generating free Cauchy system. After that quotient its Green pairings are nondegenerate (indefinite is allowed). The standard smooth time-slice argument represents every smooth Cauchy datum by applying the free equation to a smoothly time-truncated solution near the slice. Thus a polynomial element central against all linear Cauchy generators is a scalar: its commutator with a linear field is exactly its Green directional derivative times \(i\hbar\), so all solution-direction derivatives vanish. This holds of graded Wick polynomials, including the anticommuting sectors. In \eqref{eq:center}, apply the argument to the first possible nonzero coupling coefficient, using the leading free term of each \(R_f(F)\) and dividing the free linear commutator by \(i\hbar\). Iterate to the retained order. There are no external odd constants in the dynamical algebra, and the difference has ghost number one; it vanishes. This proves \eqref{eq:keycharge} with \(\Cq=\Cq^{\rm N}\). The on-shell center cannot be replaced by an assertion of equality of off-shell currents: even the free current divergence contains the free Euler products.

The bulk BV differential is nilpotent on the local dynamical quotient by the normalized plateau QME. By the graded Jacobi identity \(\Cq\star\Cq\) commutes with this same generating algebra. Its ghost number two removes its possible central term by the same coefficient argument. In the real normalization the adjoint condition for the charge follows in the same way from that of its derivation (with the fixed Krein convention in the graded sector). All these operations are formal Wick operations on the retained degrees.

For two unit temporal smearings in the common plateau their difference is \(\dd\beta\) with compact primitive, so conservation of \(j_{\rm Q}\) gives equality of their charges. One can also present the same causal algebra using auxiliary cutoffs agreeing on an open comparison strip containing a full Cauchy slice. If the alteration \(W=W_-+W_+\) is respectively supported in the past and future of that strip, set \(U_-=S(V)^{-1}\star S(V+W_-)\). For a test \(F\) in the strip causal factorization gives
\[
 S_{V+W}(F)=U_-^{-1}\star S_V(F)\star U_- .
\]
The future piece cancels in the relative matrix, the past piece enters by conjugation. Differentiation in the current and observable slots identifies the same \(\Cq\), not \(q_0\) and \(R_f(J_f^-)\) separately. The field-equation time-slice property and these embeddings yield the same physical quotient maps with composition, or one may simply use a single \(V_f\) for every artificial cut. This is the statement of Theorem~\ref{thm:charge}; it supplies no equality between separately normalized boundary theories.


\section{One graded realization and a shared physical preparation}\label{sec:state}

\subsection{An exact compatible extension of the reduced covariance}

Use the trace-reversed de Donder Einstein--Faddeev--Popov linearization. Its physical Jacobi operator \(P_{\rm E}\) has gauge map \(K\) (trace-reversed symmetrized gradient); \(D=P_{\rm E}+KK^*\) and \(R=K^*K\) are the normally hyperbolic tensor and vector operators in this normalization, and \(P_{\rm E}K=0\). Let \(E_2,E_1\) be their Cauchy data spaces at one compact Milne slice and \(q_2,q_1\) their conserved nondegenerate Hermitian Green forms. Write \(\mathcal K:E_1\to E_2\) for the induced gauge map; \(^{\dagger}\) denotes the adjoint relative to these (generally indefinite) forms. The Green identities identify \(\mathcal K^\dagger\) with the induced constraint and give \(\mathcal K^\dagger\mathcal K=0\). The reduced phase space is \(\ker\mathcal K^\dagger/\operatorname{ran}\mathcal K\).

Here \(\mathcal K\) is injective. If \(K\xi=0\), its invertible trace reversal implies that \(\xi\) is Killing. A Killing lapse on the CMC slice obeys \((-\Delta_q+3/t^2)u=0\), since it preserves the constant mean curvature, and vanishes by compactness. The tangential part would then be a spatial Killing field, excluded by the Bochner integral for negative Ricci curvature. The principal symbol of the Cauchy map is injective on either characteristic sign at nonzero spatial covector: in each sign it is the trace reversal of the symmetrized product of a nonzero null covector with the gauge amplitude. That product vanishes only when the gauge amplitude is zero. Thus \(\mathcal K\) is left elliptic in balanced Cauchy Sobolev orders.

The compact Einstein result of G\'erard provides positive physical Hadamard covariances; its microlocal split gives Cauchy projections \(c_j^\pm\) self-adjoint for \(q_j\) with
\(
 c_2^\pm\mathcal K-\mathcal K c_1^\pm\in\Psi^{-\infty}
\)
\cite[Thm.~1.1, Lemma~3.2]{Gerard}. To turn this into an \emph{exact} Ward-compatible full free covariance, take positive auxiliary Cauchy Sobolev scalar products in which the two complementary \(c_j^\pm\) are orthogonal. One may use the sum of the positive balanced norms of \(c_j^+u\) and \(c_j^-u\), symmetrized under real conjugation. This choice does not change \(q_j\). For its positive Hilbert adjoint \(^{\sharp}\), ellipticity and injectivity give the pseudodifferential left inverse
\[
 L=(\mathcal K^\sharp\mathcal K)^{-1}\mathcal K^\sharp .
\]
The smoothing intertwining error, its Hilbert adjoint, and inversion show that \(L\) intertwines signs modulo smoothing. So does \(B=L^\dagger\), since the \(c_j^\pm\) are Green-self-adjoint. Now set
\begin{equation}
 J=B-\tfrac12\mathcal K B^\dagger B,\qquad
 \Pi=1-\mathcal K J^\dagger-J\mathcal K^\dagger .
 \label{eq:nullsplit}
\end{equation}
From \(L\mathcal K=1\) and \(\mathcal K^\dagger\mathcal K=0\) we get
\(\mathcal K^\dagger J=J^\dagger\mathcal K=1\) and \(J^\dagger J=0\). Thus \(\Pi\) is an exact Green-orthogonal projection, and
\[
     E_2=\Pi E_2\oplus\mathcal K E_1\oplus J E_1 .
\]
The last two summands are isotropic and paired by \(q_1\); \(\Pi E_2\) represents the physical quotient. All these splittings intertwine the appropriate microlocal signs modulo smoothing, and can be chosen real.

Transport the positive reduced covariance to \(\Pi E_2\), denoting its operators by \(C_{\rm ph}^\pm\) with two-point forms \(\pm q_{\rm ph}C_{\rm ph}^\pm\). With \(c_1^\pm\) for the vector sector, put
\begin{equation}
      C_2^\pm=C_{\rm ph}^\pm\Pi+
                \mathcal K c_1^\pm J^\dagger+
                J c_1^\pm\mathcal K^\dagger .
      \label{eq:covariance}
\end{equation}
The metric two-point forms are \(\pm q_2 C_2^\pm\). Green adjoints, their full commutator, and the identity \(C_2^\pm\mathcal K=\mathcal K c_1^\pm\) follow exactly from \eqref{eq:nullsplit}. To see the Hadamard property, compare the physical block to \(\Pi c_2^\pm\Pi\). The latter has correct signs modulo smoothing and the same physical commutator; the difference in the plus forms equals the minus difference, so its propagated wavefront is contained in both disjoint sign relations and is empty. The null-pair terms in \eqref{eq:covariance} match the corresponding compressions of the reference modulo smoothing by the sign intertwiners. Propagate by the Cauchy Green maps. Now \cite[Prop.~3.4, Eqs.~(3.27)--(3.28)]{WZ} gives the full metric, multiplier and ghost/antighost covariance, including the mixed terms, with the exact total commutator and free BRST compatibility. There is no positive norm asserted on the ghost sector.

A chosen finite smooth real symplectic span of packet pairs may be given an exact compatible oscillator vacuum while leaving the rest Hadamard: decompose its physical symplectic orthogonal complement, restrict the original positive covariance to that complement, use the pure finite oscillator block on the span, and set the cross terms to zero. This preserves the commutator and positivity, and changes the two-point distribution only by smooth finite-rank terms. The smoothness follows from the smooth selected Cauchy data and their images under the Hadamard Cauchy operators. In particular purity of the initial infinite-dimensional reference is not assumed. The finite change is made in \(C_{\rm ph}^\pm\) in \eqref{eq:covariance}, without changing the exact null-pair construction.

\subsection{Finite-order domain and one physical state}

In the free graded realization, the null-pair and Faddeev--Popov blocks have the standard quartet form. In the matched \(E_1\) coordinates the BRST arrows and their elementary contracting maps act between the paired copies. For example one chooses an auxiliary Hilbert majorant from the absolute value of the nondegenerate \(q_1\) blocks on the two microlocal signs, in the balanced orders; identical choices on the pairs give the same smooth contraction on each copy. Taking the graded sum on the finite-particle tensor powers and dividing by the quartet number on its positive sectors gives a free homotopy \(H_0\) with
\begin{equation}
     q_0 H_0+H_0 q_0=1-\iota P .
     \label{eq:homotopy}
\end{equation}
Here \(P\) projects to the free physical sector, of ghost number zero, and the normalization of \(H_0\) follows that of free charge left multiplication. The one-particle maps in the quartet preserve smooth Cauchy data and intertwine microlocal signs; any doubled realization for a mixed physical covariance needs no contraction on its physical summands.

We specify the domain used for the retained Wick coefficients. Free smooth data admit compactly time-supported representatives: if \(S\) is free solution propagation, \(\rho\) its Cauchy trace, \(E_{\rm c}\) the causal propagator, and \(\zeta\) is a temporal step, set \(j(u)=[P_{\rm kin},\zeta]Su\) for the relevant normally hyperbolic kinetic block. Then \(\rho E_{\rm c}j(u)=u\), with the common sign convention for \(E_{\rm c}\). Eliminate the multiplier first by its algebraic equation; it can then be recovered by its differential expression. A free Cauchy map \(M\) in the split is represented by \(jM\rho E_{\rm c}\). The two Cauchy propagation relations are null Fourier integral relations preserving the chosen characteristic sign, and the spatial pseudodifferential maps in \eqref{eq:nullsplit} have diagonal Cauchy relation. Their compositions, with the smooth compact time cutoff, have no one-sided null leg with a zero cotangent on the other side. Smoothing errors have no surviving such relation. For products the absence of a common future or past orientation includes tuples with zero entries, other than the simultaneous zero.

At the finitely many retained orders the ordinary local Epstein--Glaser interacting charge coefficients are compact polynomial Wick kernels with the microcausal exclusion. Acting on smooth finite-particle coefficients gives smooth output Cauchy coefficients within each tested finite degree. Explicitly, contract the annihilation variables against smooth propagated inputs and push the creation legs to their Hadamard Cauchy sign. A singular output would require a forbidden all-oriented nonzero tuple of the original kernel with zeros in the already smoothly tested legs. Compactness of the temporal support and of the spatial section ensures proper pushforward. The same reasoning allows the sign-preserving free maps just described on these coefficient vectors. It does not establish an interacting unbounded operator domain at fixed nonzero coupling or a general renormalized time product of nonlocal insertions.

One may therefore lift a free physical vector on the retained smooth domain using the actual charge of Theorem~\ref{thm:charge}. Write \(\Cq=q_0+\sum_{m\geq1}\kappa^m Q_m\). With \(\Omega_0\) the chosen physical free state vector, define coefficientwise
\[
   \Omega_n=-H_0\sum_{m=1}^n Q_m\Omega_{n-m}.
\]
The argument of the homotopy is free-closed by nilpotence of \(\Cq\), and has zero free physical projection by ghost number. Equation~\eqref{eq:homotopy} proves \(\Cq\Omega=0\) to the retained order; the formal normalization fixes the leading free norm. For the finite test the pairing is the induced positive free physical pairing at leading nonzero order, in agreement with the formal BRST positivity deformation \cite[Thm.~4]{DF}. One does not assume the bare free cyclic functional is already invariant under the interacting physical generator. Nor does covariance under \(R_f^{-1}Q_0R_f\) suffice for this purpose: the sourced commutator in \eqref{eq:RS} must first be identified and subtracted.

Restrict this \emph{same} global state through any of the causal embeddings at the end of Section~\ref{sec:ward}. The on-shell charge, state pairing, and the external observable classes agree under these maps. On cyclic classes the induced map is the identification of the same polynomial applied to the same vector in the common realization. It intertwines the differential and preserves its radical quotient, hence preserves the physical pairing. In particular successive recuts compose without a change of state or a new projection at the middle cut.


\section{The finite graviton benchmark}\label{sec:benchmark}

\subsection{A fixed geometric canonical readout}

Use a single local constant-mean-curvature (CMC) branch on compact Milne, with its reference hyperbolic marking, and the near-identity harmonic map to fix the spatial identification. The CMC graph derivative is \(-\Delta_q+3/t^2\), since \(K=-q/t\) on the background and the slice is vacuum. It is invertible. The negative-Ricci harmonic-map Jacobi operator is likewise invertible at the reference. The effective time label in the benchmark is the CMC label reparametrized by background Milne proper time. These local elliptic gauge inverses select the same regular branch for preparation, dynamics and both exterior readouts; we assume no global CMC foliation beyond this patch.

For compactly supported diffeomorphisms near the identity the marked geometric Cauchy data in this gauge are invariant, since both elliptic gauge equations transform naturally on their domain with the target marking fixed. Their first metric and momentum coordinates are the physical TT variables. Use the exponential trace-free tensor coordinate \(q=t^2\gamma\exp(\kappa h)\) on the tensor tangent, extending by the orthogonal York scalar/longitudinal coordinates. Use its Legendre-dual momentum from the pullback of the matched ADM boundary potential \(\int\pi^{ij}\delta q_{ij}\); the intrinsic quadratic terms of these canonical coordinate functions are retained in the endpoint marks.

The first gauge compensators vanish on physical linear TT data. A second-order change of slice or spatial coordinates, acting on the reference itself, contributes a symmetrized gradient, a normal Hessian or a pure trace in the metric/momentum pair. A reference TT test annihilates all such terms by integration by parts on the compact slice. This is not a justification for dropping the quadratic terms intrinsic to the chosen canonical chart. It gives a simple form for each first nonlinear packet mark: its corrections use the first elliptic gauge inverses, differential operators, and multiplication by smooth background coefficients and packet coefficients. After pairing with a smooth packet, terms with two elliptic legs are spatial bilinear kernels of finite compositions of such operators and their transposes. At a smoothly sampled exterior time these quadratic kernels are finite sums with the form
\[
   w(t)\,\delta(t-t')\,A_t(x,x')
\]
and finite derivatives, where \(A_t\) is a smooth family of spatial pseudodifferential kernels (smoothing allowed). Their temporal wavefront components are opposite. A nonzero tuple therefore cannot have all entries future-causal or all past-causal, even if a spatial component vanishes. When the spatial part is singular on the slice diagonal its components are also opposite. These kernels are compact microcausal Wick test kernels, including spatially nonlocal terms; no sharp-time nonlinear Wick restriction is made.

The leading local gauge law and the geometric first dressing give \(s_0F_1+s_1F_0=0\) on shell, by expanding invariance under the domain diffeomorphisms. The BRST-compatible Wick convention in Section~\ref{sec:state} implements the free part exactly on these Wick coefficients. At this tree order the retarded product of an interaction with \(F_0\) has just one contraction into a smooth linear packet; \(F_1\) is only a free Wick insertion. Thus the first-order marks represent physical classes in the tree-order truncation of the common quotient. No time-ordered product of a higher-order elliptic marker with another interaction, or quantum completion of the marker to all orders, is needed for the first nonzero norm-square.

The classical Peierls--canonical identification needed for the retarded tree follows from the Green identity for the gauge-fixed Jacobi operator: for compact sources the Peierls pairing is the conserved symplectic pairing of their causal variations. The Einstein variational boundary form on the chosen compact cut is the matched ADM symplectic form; on the physical quotient the gauge doublets drop out. Solving the same field equations and then pulling to the unique near-reference marked gauge is therefore equivalent to reduced canonical evolution for the designated readings. At first order in the interaction this equivalence is also the tree part of the retarded single contraction. It does not equate actions with different canonical endpoint observables.

\subsection{An exact compact tensor triad with nonzero Einstein cubic}

The Einstein vertex used here is the true reduced bulk cubic at this background. The linear scalar and momentum constraints vanish on exact TT metric and mixed extrinsic-curvature perturbations: the background metric and Ricci endomorphism are isotropic, scalar curvature varies by a double divergence and trace on this sector, and the momentum variation is its divergence. Consequently lapse/shift corrections starting at second order only multiply the first constraints at the cubic level. The quadratic variational form has no mixed TT/complement terms on these compact constant-curvature slices. To see this, split the metric and momentum into the reference TT sector, trace parts and symmetrized gradients, integrate by parts on the closed slice, and use the constant-curvature commutators for the spatial second-order operators, which preserve TT. Time differentiation in the fixed reference chart also preserves TT. The background solves the bulk equations, so the first action variation has only the matched boundary terms. Thus solving the complementary variables at second order introduces no further bulk cubic from the quadratic form. The boundary potential fixes the canonical endpoint normalization. In the exponential metric the only possible tensor kinetic cubic with two time derivatives has factor \(\tr(\dot h[\dot h,h])=0\), since
\(e^{-\kappa h}\partial_t e^{\kappa h}=\kappa\dot h+
(\kappa^2/2)[\dot h,h]+O(\kappa^3)\).
The pullback to Legendre-dual momentum yields the same absence of a principal \(h\pi^2\) cubic; quadratic trace/longitudinal fixes pair to zero with a reference TT momentum, and the \(D\exp\) variation is the same trace commutator. The nonzero background extrinsic curvature and spatial curvature contribute only lower-derivative symbols here.

At a fixed time, up to a nonzero common Einstein factor (including the Hamiltonian sign), the spatial principal cubic in the exponential coordinate is
\begin{equation}
 \int_\Sigma (2h_{ik}h_{jl}-h_{ij}h_{kl})
                           D_kD_l h_{ij}\,\dd v_q .
 \label{eq:vertex}
\end{equation}
This is the expansion of the spatial intrinsic curvature density at unit determinant after integrating spatial divergences; the flat two-derivative polynomial is, for example, the Einstein tensor action contribution in \cite[Eq.~(3.1)]{Cabass}. Moving to the hyperbolic slice changes the connection, volume and lower symbols but not the normal-coordinate principal symbol. No temporal integration by parts with unmatched readouts is used.

Here is a finite nonzero test of \eqref{eq:vertex}. At the center of a spatial normal chart choose covectors \(k_1=(1,0,0)\), \(k_2=(0,1,0)\), \(k_3=(-1,-1,0)\). Choose their transverse trace-free polarizations
\[
 P_1=\operatorname{diag}(0,1,-1),\quad
 P_2=\operatorname{diag}(1,0,-1),\quad
 P_3=v\otimes v-e_z\otimes e_z,\quad v=\frac{e_x-e_y}{\sqrt2}.
\]
The ordered coefficient in \eqref{eq:vertex} for entries \(a,b,c\) is
\[
 -\bigl(2(P_a)_{ik}(P_b)_{jl}-(P_a)_{ij}(P_b)_{kl}\bigr)
                         k_{c,k}k_{c,l}(P_c)_{ij}.
\]
For \(123,132,213,231,312,321\) these six values are \(5/2,1/2,5/2,1/2,3/2,3/2\), so the symmetrized polynomial coefficient is \(9\). For instance the two terms can be written as
\(-2(P_ak_c)^{\mathsf T}P_c(P_bk_c)+
 \tr(P_aP_c)\,k_c^{\mathsf T}P_b k_c\), allowing direct verification. These are spatial data for a finite-time experiment, not an assertion of energy conservation of an infinite-time three-graviton scattering amplitude.

To produce exact smooth compact TT data, use the derivative \(\mathsf C'\) of the Cotton--York tensor on \((\Sigma,\gamma)\) \cite{Beig}. Cotton--York vanishes on this conformally flat background and is identically symmetric, trace-free and divergence-free in dimension three. Its derivative therefore lands in exact TT tensors. On the TT polarization symbol at nonzero spatial covector it is a nonzero constant times the second-order scalar principal symbol of the Schouten derivative followed by symmetrized curl. In a frame with \(k\) on the third axis the remaining two entries are a traceless symmetric \(2\times 2\) matrix; multiplication by the alternating tensor rotates them with complex eigenvalues of nonzero modulus proportional to \(|k|\). Thus the order-three Cotton symbol is invertible on the two-dimensional physical symbol space for every \(k\ne0\).

Use linear coordinate phases \(\phi_1=x\), \(\phi_2=y\), \(\phi_3=-x-y\) in the normal chart and a nonnegative smooth compact envelope \(a_0\). Extend each polarization smoothly over a small chart in the TT principal subspace for \(d\phi_j\), and choose a smooth seed tensor \(u_j\) with Cotton principal image \(a_0P_j(x)\). For finite large carrier \(\lambda\) set
\begin{equation}
        h_j=\lambda^{-3}\mathsf C'(u_j e^{i\lambda\phi_j}) .
        \label{eq:cotton}
\end{equation}
These are exactly TT and compactly supported. After factoring out the phase their leading amplitude is \(a_0P_j\) with a smooth expansion in \(\lambda^{-1}\). Choose mixed extrinsic-curvature or equivalently normalized canonical TT momentum on each leg of the form \(-i\omega_jh_j\) with \(\omega_j\) positive of order \(\lambda\); for the physical incoming mode use the conjugate spatial phase of the third leg and choose its positive polarization frequency. All legs evolve by the actual free equations. At the chosen slice the phases in the coefficient between the incoming mode and the two outgoing conjugates multiply to one (the triple of signs in the symbol may be reversed together). Continuity and a sufficiently small common support give a real principal integral with a strict positive sign in the convention of the displayed symbol. Thus it is \(\lambda^2 c_0+O(\lambda)\) with \(c_0\ne0\), before nonzero individual mode normalizations. Connection, Cotton lower-order and envelope terms give the indicated lower order; no long-time high-frequency estimate is needed.

The projections of these Cauchy packets onto one another and onto conjugate packets at the initial slice are rapidly small after scaling, by nonstationary phase and compact supports. Each chosen pair has nonzero symplectic normalization because \(\omega_j>0\) with the positive-frequency sign. A finite-dimensional symplectic Gram adjustment makes the three smooth pairs mutually canonical and changes the local leading coefficient by a rapidly small error. They have the exact finite-mode vacuum of Section~\ref{sec:state}. One fixes a sufficiently large \emph{finite} \(\lambda\), relative to background and envelope scales. If the result is interpreted within an effective theory the scale hierarchy can be maintained, for example by taking a sufficiently large Milne radius. The proof does not require a carrier beyond the regime of the fixed Einstein tree description.

\subsection{Unswitched evolution, finite samplers, and interference}

Let \(a,b,c\) now denote the normalized incoming and two outgoing packet annihilators in that common free convention, with \(a\) orthogonal to both \(b\) and \(c\), and \(b\perp c\). Backtransport the free labels in both exterior samplers to the same canonical reference. Choose smooth nonnegative unit time densities \(w_a,w_b\) with disjoint ordered supports inside the plateau, each of fixed positive width; the subscript on a density is an exterior-window label. Let \(B_b\) be the product of the two outgoing annihilator marks sampled in the later window, and \(B_a\) the corresponding early comparison, with the same free labels. Their zeroth-order parts both equal \(bc\). We use a product of the sampled one-packet marks; at first order the product's response is the sum of the individual responses on the same profile. Thus in the connected coefficient of
\[
             \omega_{\rm ph}\bigl((B_b-B_a)a^\dagger_{\rm in}\bigr)
\]
the first correction of the one chosen state and the first correction of the incoming mark multiply zero. This is a relative readout in one theory, not the assertion that interacting vacua at both times are the same Gaussian.

The unswitched retarded equation for a fixed interaction-picture canonical label has first coefficient the past integral of its commutator with the interaction-picture cubic \(H_3^{I}(s)\). Averaging the observation time against \(w_r\) inserts \(\int_s^\infty w_r(t)\dd t\) at time \(s\). Subtracting early from late gives the nonnegative difference
\begin{equation}
 \chi_{ab}(s)=\int_{-\infty}^s(w_a(t)-w_b(t))\,\dd t .
 \label{eq:window}
\end{equation}
It is one between the two supports and zero outside their common exterior interval. This is a retarded response to the selected external measurements, not an inserted switch of the gravitational interaction. At the retained order the normalization and tree identification in Section~\ref{sec:benchmark} give the amplitude
\begin{equation}
 -\frac{i\kappa}{\hbar}
   \left\langle bc\left|\int\chi_{ab}(s)H_3^{I}(s)\,\dd s\right|a\right\rangle .
 \label{eq:amplitude}
\end{equation}
Indeed the term with \(bc\) acting on a one-particle incoming vector is zero, leaving the indicated commutator component. A linear normal-order contraction could only leave the incoming packet as a spectator after creating one additional packet; the two tested outgoing modes are orthogonal to it, so this and disconnected contractions vanish. The vacuum Wick normalization of the three selected free modes is exact, and unobserved free modes may be in a mixed reference covariance.

With \(\sqrt{\hbar}\) per canonically normalized external quantum leg, write \(f_3(s)=\hbar^{-3/2}\langle bc|H_3^{I}(s)|a\rangle\). It is smooth and nonzero at the slice of the tensor calculation. There is therefore a constant phase \(\theta\) and a short open interval in which \(\Re(e^{i\theta}f_3(s))\) is strictly positive. Place both fixed-width samplers and their gap inside that interval. Then
\[
       \mathfrak a=\int\chi_{ab}(s)f_3(s)\,\dd s\ne0 .
\]
This proves \eqref{eq:cubicform}; no sharp-width limit, asymptotic rate, or noncollinear infinite-time scattering is required. In particular, the cancellation of first-order incoming and state corrections is sufficient here. The nonlinear brackets of averaged coordinates need not equal those of the central time coordinates, and we do not use such an equality.

For the retained positive contribution take the physical pairing on the normalized outgoing two-packet vector. The number is the squared leading matrix element, \(\hbar\kappa^2|\mathfrak a|^2\), or equivalently the evaluation of the leading transition followed by the finite-rank projection in this tested physical representation. A rank-one projection in the local field algebra is not needed. Both legs of this selected intermediate pairing are smooth finite-rank Wightman contractions. With zero zeroth-order overlap, second-order changes of the amplitudes or local first-loop contacts in an unrelated elastic amplitude do not change this leading square. This does not evaluate or impose a normalization on all other one-loop graphs.

If \(c\) labels an internal sampling profile, the elementary identity
\(\chi_{ac}+\chi_{cb}=\chi_{ab}\) yields \(\mathfrak a=\mathfrak a_{ac}+\mathfrak a_{cb}\). Hence a middle bookkeeping cut produces, in the same physical pairing,
\[
 |\mathfrak a|^2=|\mathfrak a_{ac}|^2+|\mathfrak a_{cb}|^2+
                     2\Re\bigl(\overline{\mathfrak a_{ac}}\mathfrak a_{cb}\bigr).
\]
For example, all three ordered samplers can be concentrated in a smaller fixed-phase interval where both partial integrals lie in a sector of width strictly less than \(\pi/2\); the cross term is then strictly positive. The equality of charge and one-state transport in Sections~\ref{sec:ward}--\ref{sec:state} explains why there is no projection or reset at an imagined intermediate cut. This verifies Theorem~\ref{thm:experiment} and distinguishes coherent causal recutting from a protocol with an actual intermediate measurement.

\section{Scope and outlook}

The construction identifies a source-dependent quantum Ward correction with an interior charge only on the dynamical quotient, under a fixed local anomaly-normalized pure Einstein prescription. The finite radiative witness additionally uses compact hyperbolic Milne, chosen geometric endpoint data, and smooth finite-order Wick tests. It proves a same-theory, cut-independent positive process, not a nonperturbative quantum gravity Hilbert space, independently defined boundary-face constraints, a complete elastic loop calculation, or all higher products of nonlocal CMC readouts.


\begin{thebibliography}{12}
\bibitem{BFR}
R.~Brunetti, K.~Fredenhagen and K.~Rejzner,
Quantum gravity from the point of view of locally covariant quantum field theory,
\emph{Commun. Math. Phys.} \textbf{345} (2016), 741--779,
\href{https://arxiv.org/abs/1306.1058}{arXiv:1306.1058}.
\bibitem{Rejzner}
K.~Rejzner, Remarks on local symmetry invariance in perturbative algebraic quantum field theory,
\href{https://arxiv.org/abs/1301.7037}{arXiv:1301.7037}.
\bibitem{RS}
K.~Rejzner and M.~Schiavina, Perturbative algebraic quantum field theory with smoothened boundary,
\href{https://arxiv.org/abs/2607.13765}{arXiv:2607.13765}.
\bibitem{Hollands}
S.~Hollands, Renormalized quantum Yang--Mills fields in curved spacetime,
\emph{Rev. Math. Phys.} \textbf{20} (2008), 1033--1172,
\href{https://arxiv.org/abs/0705.3340}{arXiv:0705.3340}.
\bibitem{TT}
M.~Taslimi Tehrani, Quantum BRST charge in gauge theories in curved space-time,
\href{https://arxiv.org/abs/1703.04148}{arXiv:1703.04148}.
\bibitem{Gerard}
C.~G\'erard, Hadamard states for linearized gravity on spacetimes with compact Cauchy surfaces,
\emph{Commun. Math. Phys.} \textbf{405} (2024), 145,
\href{https://arxiv.org/abs/2311.11043}{arXiv:2311.11043}.
\bibitem{WZ}
M.~Wrochna and J.~Zahn, Classical phase space and Hadamard states in the BRST formalism for gauge field theories on curved spacetime,
\emph{Rev. Math. Phys.} \textbf{29} (2017), 1750014,
\href{https://arxiv.org/abs/1407.8079}{arXiv:1407.8079}.
\bibitem{DF}
M.~D\"utsch and K.~Fredenhagen, A local (perturbative) construction of observables in gauge theories: the example of QED,
\emph{Commun. Math. Phys.} \textbf{203} (1999), 71--105,
\href{https://arxiv.org/abs/hep-th/9807078}{arXiv:hep-th/9807078}.
\bibitem{Cabass}
G.~Cabass, D.~Stefanyszyn, J.~Supe\l{} and A.~Thavanesan,
On graviton non-Gaussianities in the effective field theory of inflation,
\href{https://arxiv.org/abs/2209.00677}{arXiv:2209.00677}.
\bibitem{Beig}
R.~Beig, TT-tensors and conformally flat structures on \(3\)-manifolds,
\href{https://arxiv.org/abs/gr-qc/9606055}{arXiv:gr-qc/9606055}.
\bibitem{P19}
\emph{A Local Einstein--Scalar BV/BRST Benchmark for Boundary Reconstruction Programs},
Research manuscript 19,
\href{https://github.com/g3173028019-beep/ai-automated-quantum-gravity-research/tree/main/papers/19-local-gauge-benchmarks}{project archive}.
\bibitem{P53}
\emph{Coherent Sewing of a Physical Gravitational Shear: Common Cut Observables, a Compact Revival, and Infinite Residual Covariance},
Research manuscript 53,
\href{https://github.com/g3173028019-beep/ai-automated-quantum-gravity-research/tree/main/papers/53-coherent-quantum-shear-sewing}{project archive}.
\end{thebibliography}


\end{document}

